\documentclass[10pt]{article}
\usepackage[T1]{fontenc}
\usepackage[utf8]{inputenc}
\usepackage[french]{babel}
\usepackage{amsmath,amssymb,mathrsfs,mathpazo}
\usepackage{geometry}
\geometry{a4paper,margin=2cm}
\usepackage{graphicx,array,tabularx,tabularray,colortbl,multirow,multicol,booktabs,enumitem,pifont,ulem,fancybox,diagbox,cancel}
\usepackage{pstricks,pstricks-add,pst-plot,pst-tree,pst-node,pst-text,pst-func,pst-math,pst-bspline,pst-3dplot,pst-eucl,pst-solides3d}
\let\pscalculate\undefined
\usepackage{pst-bezier}
\usepackage{tikz,pgfplots,tkz-tab}
\usetikzlibrary{arrows,arrows.meta,patterns,calc,positioning,trees,decorations.markings,shapes.geometric,shapes.misc,shapes.arrows,intersections,decorations.pathmorphing,decorations.pathreplacing,backgrounds,shadows,babel}
\pgfplotsset{compat=1.18}
\usepgfplotslibrary{fillbetween}
\usepackage[np]{numprint}
\usepackage{listings,makecell,pas-tableur,twoopt}
\let\barre\undefined
\usepackage[pstricks]{tablvar}
\definecolor{darkgreen}{RGB}{0,100,0}
\providecommand{\red}{\color{red}}
\providecommand{\blue}{\color{blue}}
\providecommand{\green}{\color{green}}
\providecommand{\black}{\color{black}}
\providecommand{\white}{\color{white}}
\providecommand{\gray}{\color{gray}}
\providecommand{\cyan}{\color{cyan}}
\providecommand{\magenta}{\color{magenta}}
\providecommand{\yellow}{\color{yellow}}
\providecommand{\violet}{\color{violet}}
\providecommand{\orange}{\color{orange}}
\providecommand{\euro}{\texteuro}
\providecommand{\EUR}{\texteuro}
\providecommand{\ieme}{\textsuperscript{e}}
\providecommand{\degres}{\ensuremath{^\circ}}
\providecommand{\degree}{\ensuremath{^\circ}}
\providecommand{\vv}[1]{\overrightarrow{#1}}
\providecommand{\Coord}[1]{#1}
\providecommand{\decosix}{}
\providecommand{\starredbullet}{\ensuremath{\bullet}}
\pagestyle{empty}
\setlength{\parindent}{0pt}
\usepackage{adjustbox}\DeclareUnicodeCharacter{2236}{:}\DeclareUnicodeCharacter{2190}{\ensuremath{\leftarrow}}
\providecommand{\R}{\mathbb{R}}
\providecommand{\N}{\mathbb{N}}
\providecommand{\Z}{\mathbb{Z}}
\providecommand{\Q}{\mathbb{Q}}
\providecommand{\D}{\mathbb{D}}
\providecommand{\C}{\mathbb{C}}
\providecommand{\e}{{\mathrm{e}}}
\providecommand{\vect}[1]{\overrightarrow{#1}}
\providecommand{\vectt}[1]{\overrightarrow{\mathrm{#1}}}
\providecommand{\vvt}[1]{\overrightarrow{\mathrm{#1}}}
\providecommand{\np}[2][]{{#2~#1}}
\providecommand{\vv}[1]{{\overrightarrow{#1}}}
\renewcommand{\d}{{\mathrm{d}}}
\providecommand{\degres}{{^\circ}}
\providecommand{\degre}{{^\circ}}
\providecommand{\blue}{\color{blue}}
\providecommand{\red}{\color{red}}
\providecommand{\enskip}{\quad}
\providecommand{\llbracket}{\lbrack\!\lbrack}
\providecommand{\rrbracket}{\rbrack\!\rbrack}
\providecommand{\euro}{{\text{€}}}
\providecommand{\no}{{\text{n°}}}
\providecommand{\interval}[2]{{[#1;#2]}}
\providecommand{\intervFO}[2]{{[#1;#2[}}
\providecommand{\dotfill}{\ldots}

\providecommand{\ds}{\displaystyle}
\providecommand{\textding}[1]{#1}
\providecommand{\Oij}{\(\mathrm{O};\vec{\imath},\vec{\jmath}\)}
\providecommand{\Oijk}{\(\mathrm{O};\vec{\imath},\vec{\jmath},\vec{k}\)}
\providecommand{\Ouv}{\(\mathrm{O};\vec{u},\vec{v}\)}
\providecommand{\intervFF}[2]{\ensuremath{[#1;#2]}}
\renewcommand{\euro}{\ifmmode\text{\texteuro}\else\texteuro\fi}
\DeclareUnicodeCharacter{00B0}{\ifmmode^{\circ}\else\textdegree\fi}
\DeclareUnicodeCharacter{20AC}{\ifmmode\text{\texteuro}\else\texteuro\fi}
% Énoncés : Éducation nationale. Transcription et corrigés : APMEP, Denis Vergès et contributeurs. Découpage et classement : CoopMaths, Sébastien Lozano, Éric Elter et contributeurs. Adaptation de la présentation : Maths Explorer.
% CC BY-SA : https://coopmaths.fr/www/about/
% https://coopmaths.fr/alea/static/bac/2023/tex/bac_2023_08_sujet2_caledonie_2_cor.tex
\begin{document}
\noindent\textit{APMEP / CoopMaths. CC BY-SA. Présentation adaptée par Maths Explorer.}\par\medskip
\medskip

On considère la fonction $f$ définie sur $[0~;~ +\infty[$ par: $f(x) = x\e^{-x}$.









\begin{enumerate}
\item 



$f(x)=x \e^{-x}=\dfrac{x}{\e^{x}}$.
Or $\ds\lim_{x\to +\infty} \dfrac{\e^{x}}{x}=+\infty$; 
donc $\ds\lim_{x\to +\infty} \dfrac{x}{\e^{x}} =0$ et donc  $\ds\lim_{x\to +\infty} f(x)=0$.

On en déduit que la courbe $\mathcal{C}_f$ possède une asymptote en $+ \infty$ d'équation $y=0$; c'est l'axe des abscisses.

\item Pour tout réel $x$ appartenant à $[0~;~ +\infty[$:
$f'(x)=1\times \e^{-x}+x\times (-1)\e^{-x}=(1-x)\e^{-x}$.

\item 

Pour tout réel $x$, $\e^{-x}>0$ donc $f'(x)$ est du signe de $(1-x)$ sur $[0~;~ +\infty[$.

On étudie le signe de $f'(x)$ sur $[0~;~ +\infty[$.

\begin{center}
{\renewcommand{\arraystretch}{1.5}
\def\esp{\hspace*{3cm}}
$\begin{array}{|c | *{5}{c} |} 
\hline
x  & 0 & \esp & 1 & \esp  & +\infty \\
\hline
1-x &  & \pmb{+} &  \vline\hspace{-2.7pt}{0} & \pmb{-} &   \\
\hline
\e^{-x} &  & \pmb{+} &  \vline\hspace{-2.7pt}{\phantom 0} & \pmb{+} &   \\
\hline
f'(x) &  & \pmb{+} &  \vline\hspace{-2.7pt}{0} & \pmb{-} &   \\
\hline
\end{array}$
}
\end{center}

$f(0)=0$; $f(1)=\e^{-1}$ et $\ds\lim_{x\to +\infty} f(x)=0$.\\
On dresse le tableau de variations de $f$ sur  $[0~;~ +\infty[$:

\begin{center}
{\renewcommand{\arraystretch}{1.3}
\psset{nodesep=3pt,arrowsize=2pt 3}
\def\esp{\hspace*{2.5cm}}
\def\hauteur{0pt}
$\begin{array}{|c|*5{c}|}
\hline
x & 0  & \esp & 1 & \esp & +\infty \\ 
\hline
f'(x) &  &   \pmb{+} & \vline\hspace{-2.7pt}0 & \pmb{-} & \\ 
\hline
 & &  &   \Rnode{max}{\e^{-1}}  &  &   \\  
f(x) & &     &  &  &  \rule{0pt}{\hauteur} \\ 
 & \Rnode{min1}{0} &   &  &  &   \Rnode{min2}{0} \rule{0pt}{\hauteur}    
 \ncline{->}{min1}{max} 
 \ncline{->}{max}{min2} 
 \\ 
\hline
\end{array} $
}
\end{center}	

\item 
$\e^{-1}\approx 0,369>\frac{367}{\np{1000}}$

\begin{list}{\textbullet}{}
\item Sur l'intervalle $[0\;;\;1]$, la fonction $f$ est continue et strictement croissante; elle va de $0$ à $\e^{-1}>0,367$. Donc d'après le corollaire du théorème des valeurs intermédiaires, l'équation $f(x)=0,367$ admet une solution unique sur cet intervalle.
\item Sur l'intervalle $[1\;;\;+\infty[$, la fonction $f$ est continue et strictement décroissante; elle va de $\e^{-1}>0,367$ à $0$. Donc d'après le corollaire du théorème des valeurs intermédiaires, l'équation $f(x)=0,367$ admet une solution unique sur cet intervalle.
\end{list}

L'équation $f(x) = \frac{367}{\np{1000}}$ admet donc deux solutions sur l'intervalle $[0~;~ +\infty[$.

\item On admet que pour tout $x$ appartenant à $[0~;~ +\infty[$ : 
$f''(x) = \e^{-x}(x - 2)$.

Pour étudier  la convexité de la fonction $f$ sur l'intervalle $[0~;~ +\infty[$, on étudie le signe de $f''(x)$.

\begin{center}
{\renewcommand{\arraystretch}{1.5}
\def\esp{\hspace*{3cm}}
$\begin{array}{|c | *{5}{c} |} 
\hline
x  & 0 & \esp & 2 & \esp  & +\infty \\
\hline
x-2 &  & \pmb{-} &  \vline\hspace{-2.7pt}{0} & \pmb{+} &   \\
\hline
\e^{-x} &  & \pmb{+} &  \vline\hspace{-2.7pt}{\phantom 0} & \pmb{+} &   \\
\hline
f''(x) &  & \pmb{-} &  \vline\hspace{-2.7pt}{0} & \pmb{+} &   \\
\hline
 &  & f \text{ est concave} &  \vline\hspace{-2.7pt}{\phantom 0} & f \text{ est convexe} &   \\
\hline
\end{array}$
}
\end{center}

On peut même préciser que la courbe admet le point d'abscisse 2 comme point d'inflexion.

\begin{minipage}{6cm}
\item Soit $a$ un réel appartenant à $[0~;~ +\infty[$ et A le
point de la courbe $\mathcal{C}_f$ d'abscisse $a$.

On note $T_a$ la tangente à $\mathcal{C}_f$ en A.

On note H$_a$ le point d'intersection de la droite $T_a$ et de l'axe des ordonnées.

On note $g(a)$ l'ordonnée de H$_a$.

La situation est représentée sur la figure ci-contre.
\end{minipage}\hfill
\begin{minipage}{7cm}
\psset{xunit=1cm,yunit=8cm,arrowsize=2pt 3}
\begin{pspicture}(-1,-0.1)(6,0.5)
\psaxes[linewidth=1.15pt,Dx=10,Dy=10]{->}(0,0)(-1,0)(6,0.5)
\psplot[plotpoints=2000,linewidth=1.25pt,linecolor=red]{0}{6}{x 2.71828 x exp div}
\uput[u](4,0.073){\red $\mathcal{C}_f$}
\psplotTangent{0.75}{1.25}{x 2.71828 x exp div}
\uput[u](0.75,0.354){A}
\psline[linestyle=dotted,linewidth=1.3pt](0.75,0)(0.75,0.354)
\psline[linewidth=0.8pt]{->}(-0.3,0)(-0.3,0.265)
\psline[linestyle=dotted,linewidth=1.3pt](0,0.265)(-0.3,0.265)
\uput[l](-0.3,0.1325){$g(a)$}
\uput[ul](0,0.265){H$_a$}
\uput[d](0.75,0){$a$}\uput[u](2,0.4){$T_a$}
\end{pspicture}
\end{minipage}

	\begin{enumerate}
		\item 
La tangente $T_a$ a pour équation:

$\begin{gathered}y=f'(a)(x-a)+f(a)
\iff
y= (1-a)\e^{-a}(x-a)+a\e^{-a}\\
\phantom{y=f'(a)(x-a)+f(a)}
\iff
y= \left [(1-a)\e^{-a}\right ]x- (1-a)\e^{-a}a+a\e^{-a}\\
\phantom{y=f'(a)(x-a)+f(a)}
\iff
y= \left [(1-a)\e^{-a}\right ]x -a\e^{-a}+a^2\e^{-a}+a\e^{-a}\\
\phantom{y=f'(a)(x-a)+f(a)}
\iff
y= \left [(1-a)\e^{-a}\right ]x +a^2\e^{-a}\end{gathered}$
		
		\item 
$g(a)$ est l'ordonnée	du point H$_a$ de la tangente d'abscisse 0;
donc 

$g(a)= \left [(1-a)\e^{-a}\right ]\times 0 +a^2\e^{-a}=a^2\e^{-a}$.
		
		\item 
Soit $g$ la fonction définie sur $[0~;~ +\infty[$ par $g(x)=x^2\e^{-x}$.

$g'(x)=2x\times \e^{-x}+ x^2\times (-1)\e^{-x}= -x(x-2)\e^{-x}$

\begin{center}
{\renewcommand{\arraystretch}{1.5}
\def\esp{\hspace*{3cm}}
$\begin{array}{|c | *{5}{c} |} 
\hline
x  & 0 & \esp & 2 & \esp  & +\infty \\
\hline
-x &  0 & \pmb{-} &  \vline\hspace{-2.7pt}{\phantom 0} & \pmb{-} &   \\
\hline
x-2 &  & \pmb{-} &  \vline\hspace{-2.7pt}{0} & \pmb{+} &   \\
\hline
\e^{-x} &  & \pmb{+} &  \vline\hspace{-2.7pt}{\phantom 0} & \pmb{+} &   \\
\hline
g'(x) &  & \pmb{+} &  \vline\hspace{-2.7pt}{0} & \pmb{-} &   \\
\hline
\end{array}$
}
\end{center}

D'après le tableau de signes de $g'(x)$, la fonction $g$ est croissante puis décroissante; elle admet un maximum pour $x=2$, ce qui correspond au point d'inflexion de la courbe $\mathcal{C}_f$.


	\end{enumerate}
\end{enumerate}

\bigskip

\end{document}
